\section{Вероятность, условная вероятность}
\label{sec:probability-and-conditional-probability}

\subsection{Случайные события и вероятностное пространство}

Пусть $\Omega$ --- множество всех возможных \emph{элементарных исходов} случайного
эксперимента. Подмножества $A\subset\Omega$, которым можно поставить в соответствие
вероятность, называются \emph{случайными событиями}. Достоверному событию соответствует
всё пространство $\Omega$, невозможному --- пустое множество $\varnothing$.

Для работы с событиями выделяют $\sigma$-алгебру $\mathcal F$ подмножеств $\Omega$.
Она содержит $\Omega$, замкнута относительно дополнения и счётных объединений
(следовательно, и счётных пересечений). Тройка
\[
    (\Omega,\mathcal F,P)
\]
называется \emph{вероятностным пространством}, где $P$ --- вероятностная мера.

Над событиями выполняются обычные операции над множествами:
\[
    A\cup B \quad\text{--- происходит хотя бы одно из событий},
\]
\[
    A\cap B \quad\text{--- происходят оба события},
    \qquad
    \overline A=\Omega\setminus A \quad\text{--- событие, противоположное $A$}.
\]
События $A$ и $B$ называются несовместными, если $A\cap B=\varnothing$.

\subsection{Аксиомы вероятности и основные свойства}

Вероятность удовлетворяет аксиомам Колмогорова:
\begin{enumerate}
    \item $P(A)\geq 0$ для любого $A\in\mathcal F$;
    \item $P(\Omega)=1$;
    \item для попарно несовместных событий $A_1,A_2,\ldots$
    \[
        P\!\left(\bigcup_{i=1}^{\infty}A_i\right)
        =\sum_{i=1}^{\infty}P(A_i).
    \]
\end{enumerate}

Отсюда следуют основные свойства:
\[
    P(\varnothing)=0,
    \qquad
    0\leq P(A)\leq 1,
    \qquad
    P(\overline A)=1-P(A).
\]
Для двух произвольных событий
\[
    P(A\cup B)=P(A)+P(B)-P(A\cap B),
\]
а для несовместных событий последний член равен нулю.

\subsection{Классическое и геометрическое определения вероятности}

Если пространство элементарных исходов конечно и все $n$ исходов равновозможны,
а событию $A$ благоприятствуют $m$ исходов, то используется классическое определение
\[
    P(A)=\frac{m}{n}.
\]

Для непрерывного множества равновозможных исходов используется геометрическая
вероятность. Если точка равномерно выбирается в области $G$, а событию $A$
соответствует попадание в подобласть $g\subset G$, то
\[
    P(A)=\frac{\operatorname{mes} g}{\operatorname{mes} G},
\]
где $\operatorname{mes}$ означает соответствующую геометрическую меру: длину,
площадь или объём.

\subsection{Условная вероятность}

Пусть $P(B)>0$. \emph{Условной вероятностью} события $A$ при условии, что событие
$B$ произошло, называется величина
\[
    P(A\mid B)=\frac{P(A\cap B)}{P(B)}.
\]
Условная вероятность показывает, как изменяется вероятность события $A$, если
известно, что произошло $B$.

Из определения непосредственно следует \emph{теорема умножения вероятностей}:
\[
    P(A\cap B)=P(B)P(A\mid B),\qquad P(B)>0.
\]
Если также $P(A)>0$, то
\[
    P(A\cap B)=P(A)P(B\mid A),
\]
поэтому при положительных $P(A)$ и $P(B)$ можно записать обе формы одновременно.
Для нескольких событий, если все возникающие условные вероятности определены,
\[
    P(A_1\cap\cdots\cap A_n)
    =P(A_1)P(A_2\mid A_1)\cdots
    P(A_n\mid A_1\cap\cdots\cap A_{n-1}).
\]

\subsection{Независимость событий}

События $A$ и $B$ называются \emph{независимыми}, если
\[
    \boxed{P(A\cap B)=P(A)P(B).}
\]
Это определение корректно и для событий нулевой вероятности. Если $P(B)>0$, оно
эквивалентно условию $P(A\mid B)=P(A)$; если $P(A)>0$ --- условию
$P(B\mid A)=P(B)$.
Для нескольких событий различают попарную независимость и независимость в
совокупности. Для независимости в совокупности требуется, чтобы для любого набора
индексов $i_1,\ldots,i_k$
\[
    P(A_{i_1}\cap\cdots\cap A_{i_k})
    =P(A_{i_1})\cdots P(A_{i_k}).
\]

Важно не путать независимость и несовместность: несовместные события с
положительными вероятностями, наоборот, зависимы, так как наступление одного
исключает наступление другого.

\subsection{Формула полной вероятности и формула Байеса}

Пусть события $B_1,\ldots,B_n$ попарно несовместны, образуют полную группу,
\[
    \bigcup_{i=1}^{n}B_i=\Omega,
\]
и $P(B_i)>0$. Тогда для любого события $A$ справедлива \emph{формула полной
вероятности}
\[
    P(A)=\sum_{i=1}^{n}P(B_i)P(A\mid B_i).
\]

Если событие $A$ уже произошло и $P(A)>0$, то вероятности гипотез $B_i$ можно пересчитать
по \emph{формуле Байеса}:
\[
    P(B_i\mid A)
    =\frac{P(B_i)P(A\mid B_i)}{P(A)}
    =\frac{P(B_i)P(A\mid B_i)}
    {\displaystyle\sum_{j=1}^{n}P(B_j)P(A\mid B_j)}.
\]
Вероятности $P(B_i)$ называют априорными, а $P(B_i\mid A)$ --- апостериорными.

\subsection{Схема Бернулли}

В схеме Бернулли производится $n$ независимых испытаний, в каждом из которых
событие $A$ происходит с одной и той же вероятностью $p$, а не происходит с
вероятностью $q=1-p$. Вероятность того, что событие $A$ произойдёт ровно $m$ раз,
равна
\[
    P_n(m)=\binom{n}{m}p^m q^{\,n-m}.
\]
Это биномиальный закон распределения вероятностей.


\subsection{Дополнительные свойства вероятностной меры}

Из аксиом Колмогорова следуют свойства монотонности и непрерывности меры. Если
$A\subset B$, то
\[
    P(A)\le P(B),\qquad P(B\setminus A)=P(B)-P(A).
\]
Если $A_n\uparrow A$, то
\[
    P(A_n)\to P(A),
\]
а если $A_n\downarrow A$ и $P(A_1)<\infty$ (для вероятности это автоматически), то
\[
    P(A_n)\to P(A).
\]
Эти свойства удобны при переходах к пределу и при работе с бесконечными
объединениями и пересечениями событий.

Для конечного семейства событий справедлива формула включений--исключений. Например,
для трёх событий
\[
\begin{aligned}
P(A\cup B\cup C)={}&P(A)+P(B)+P(C)
-P(A\cap B)-P(A\cap C)-P(B\cap C)\\
&+P(A\cap B\cap C).
\end{aligned}
\]

\subsection{Полная группа событий и вывод формул полной вероятности и Байеса}

Если $B_1,\ldots,B_n$ образуют разбиение $\Omega$, то
\[
    A=\bigcup_{i=1}^n (A\cap B_i),
\]
причём события $A\cap B_i$ попарно несовместны. Поэтому
\[
    P(A)=\sum_{i=1}^n P(A\cap B_i)
        =\sum_{i=1}^n P(B_i)P(A\mid B_i).
\]
Это и есть формула полной вероятности. Формула Байеса получается из определения
условной вероятности:
\[
P(B_i\mid A)=\frac{P(A\cap B_i)}{P(A)}
=\frac{P(B_i)P(A\mid B_i)}{\sum_j P(B_j)P(A\mid B_j)}.
\]
Таким образом, Байес не является отдельной аксиомой: это непосредственное следствие
условной вероятности и полной вероятности.

\subsection{Попарная и совместная независимость}

Для трёх и более событий попарная независимость слабее независимости в совокупности.
Например, могут выполняться
\[
P(A_i\cap A_j)=P(A_i)P(A_j)\qquad (i\ne j),
\]
но не выполняться
\[
P(A_1\cap A_2\cap A_3)=P(A_1)P(A_2)P(A_3).
\]
Поэтому при формулировке схемы Бернулли подразумевается независимость всех испытаний
в совокупности.

\subsection{Числовые характеристики числа успехов в схеме Бернулли}

Если $\nu_n$ --- число успехов в $n$ испытаниях Бернулли, то
\[
    \nu_n\sim\operatorname{Bin}(n,p),\qquad
    \mathbb E\nu_n=np,\qquad
    \operatorname{Var}\nu_n=np(1-p).
\]
Относительная частота $\nu_n/n$ сходится к $p$ по вероятности; это содержание
теоремы Бернулли как частного случая закона больших чисел.

\subsection{Что важно сказать на экзамене}

Вероятность формально задаётся вероятностным пространством
$(\Omega,\mathcal F,P)$ и аксиомами Колмогорова. Для условной вероятности нужно
обязательно записать
\[
    P(A\mid B)=\frac{P(A\cap B)}{P(B)},
\]
после чего естественно получить теорему умножения. Независимость событий
эквивалентна равенству $P(A\cap B)=P(A)P(B)$. Если имеется полная группа гипотез,
используются формула полной вероятности и формула Байеса.

\newpage
